\section*{Chapter \ref{ch:rc:bermudans-and-callables} --- Bermudans and Callables}

\begin{solution}{exo:rc:bermudans-and-callables:1}
Lower bound: the three-year European, 249.5 basis points. Upper bound: the sum of the
nine, 1\,647.5. The Bermudan, 358.4, lies between, much nearer the lower bound.
\end{solution}

\begin{solution}{exo:rc:bermudans-and-callables:2}
A European's value is roughly annuity times volatility times the square root of
expiry: the one-year has the longest swap (largest annuity) but little time for
rates to move; the nine-year has time but a one-year swap. The product peaks at
three years here.
\end{solution}

\begin{solution}{exo:rc:bermudans-and-callables:3}
She has sold the issuer a Bermudan option to redeem at the accreted value, a
receiver-type option on long rates (the issuer calls when rates fall). The issuer
passes it through the swap to the dealer, who hedges it by selling swaption
volatility to the market.
\end{solution}

\begin{solution}{exo:rc:bermudans-and-callables:4}
With fewer exercise dates left, waiting is worth less, so exercise needs a smaller
move in the holder's favour; at the last date the Bermudan is a European and
exercise happens whenever the swap is in the money. The boundary stays below the
forward because at the money the continuation value (option on later dates)
exceeds the immediate intrinsic value.
\end{solution}

\begin{solution}{exo:rc:bermudans-and-callables:5}
The best European is held fixed by recalibrating $\sigma$ on the five-into-five, so
it moves little. Stronger mean reversion makes swap rates of different exercise
dates less correlated: a low rate at year three says less about year six, so the
right to choose the date later is worth more (101.9, 109.0 and 120.0 basis points).
\end{solution}

\begin{solution}{exo:rc:bermudans-and-callables:6}
Fitted and applied on the same paths, the regression uses each path's own future
to decide whether to exercise, a look-ahead that biases the price upward; applied
to independent paths, the rule is a legitimate (sub-optimal) stopping rule, whose
expected payoff is a lower bound.
\end{solution}

\begin{solution}{exo:rc:bermudans-and-callables:7}
The callable zero is worth par at an accretion yield of 5.61\%; the straight
thirty-year zero on the curve yields 4.13\% (annual). The investor earns 1.48
percentage points a year for the call.
\end{solution}

\begin{solution}{exo:rc:bermudans-and-callables:8}
The European covers only the lower bound: the \omterm{def:rc:bermudans-and-callables:switch}{switch option} (here 44\% of the
best European) is unhedged, and the vega of the Bermudan is spread over all the
co-terminals, shifting as rates move; the notional of a delta- and vega-neutral
hedge is not the Bermudan's notional. Hedge each co-terminal's vega and delta, and
reserve for the switch.
\end{solution}

\begin{solution}{pb:rc:bermudans-and-callables:1}
\textbf{1.} 4.13\%.
\textbf{2.} 5.61\%.
\textbf{3.} 1.526.
\textbf{4.} 0.526 per unit issued: 1.526 minus par.
\textbf{5.} For 1.48 percentage points a year of extra yield, and because the regulation
made these bonds the most efficient way to hold long dollar duration.
\textbf{6.} Year five, the first call: 0.456.
\textbf{7.} 0.070 per unit.
\textbf{8.} A zero-coupon bond called at its accreted value is most in the money early
(the accretion compounds above the curve's rates from the start), and the first
call cuts the longest remaining liability.
\textbf{9.} A Bermudan receiver-type option on long rates, received through the swap
from the issuer, against which the dealer pays the floating funding.
\textbf{10.} It is a large long-volatility position in illiquid long-dated options;
dealers are paid to warehouse and hedge it, not to hold it.
\textbf{11.} USD~2.33 million per basis point of the model's $\sigma$ per USD~1 billion.
\textbf{12.} Long-dated swaptions (co-terminals of thirty-year swaps from five years
on), and longer-expiry receivers.
\textbf{13.} About USD~23.3 million per basis point of $\sigma$.
\textbf{14.} It is pushed down by the steady selling; dealers long the Bermudans mark
them lower as volatility falls.
\textbf{15.} The Bermudan disappears when called: the dealer is left short the
volatility it sold, and must buy it back, often in a falling-rate, rising-volatility
market.
\textbf{16.} Reinvestment risk when called (in low-rate states), currency risk against
its liabilities, and concentration in few issuers.
\textbf{17.} Few dealers can warehouse and hedge the options; their hedging is
correlated, so the flows move the market and a reversal (calls) hits them all at
once.
\textbf{18.} The volatility and correlation structure of the model (the switch value),
the smile of long-dated swaptions, and the exercise rule.
\textbf{19.} Named result: \emph{the Formosa hedge}: the callable zero is at par at an
accretion of 5.61\% against 4.13\% for the straight zero; the Bermudan call is worth
0.070 more than the best European (0.456); its vega is USD~2.33 million per basis
point per USD~1 billion.
\textbf{20.} The dealer sells volatility, because it bought a Bermudan it does not
want to carry, from an issuer who took it from an investor paid to give it up.
\end{solution}

\begin{solution}{iq:rc:bermudans-and-callables:1}
The right to enter, on any one of several dates, a swap ending on a fixed date. It
is worth at least any single European (commit to that date) and more, because the
holder chooses the date after seeing rates: the \omterm{def:rc:bermudans-and-callables:switch}{switch option}.
\iqlookfor{the definition and the \omterm{def:rc:bermudans-and-callables:switch}{switch option}.}
\end{solution}

\begin{solution}{iq:rc:bermudans-and-callables:2}
Simulate the state; from the last exercise date back, regress the discounted
realised future cash flows of in-the-money paths on basis functions of the state
and exercise where the exercise value exceeds the fitted continuation. Any
exercise rule gives an achievable payoff, so its value (on paths independent from
the fit) is at most the optimum.
\iqlookfor{the backward regression and why it bounds from below.}
\end{solution}

\begin{solution}{iq:rc:bermudans-and-callables:3}
Mean reversion: it decorrelates the swap rates of successive exercise dates; with
the Europeans' prices held by recalibrating $\sigma$, higher $\kappa$ gives a larger
\omterm{def:rc:bermudans-and-callables:switch}{switch option}.
\iqlookfor{mean reversion as correlation between exercises.}
\end{solution}

\begin{solution}{iq:rc:bermudans-and-callables:4}
Issuers of callable bonds swap them with dealers and pass on the call; dealers,
long Bermudans, hedge by selling long-dated swaptions, a steady supply that
depresses long-dated implied volatility; calls or unwinds reverse the flow.
\iqlookfor{the chain from investor to swaption market.}
\end{solution}

\begin{solution}{iq:rc:bermudans-and-callables:5}
Compare with Europeans (single-date limit) and bounds; tree against regression on
the same model; convergence in steps and paths; a dual upper bound for the gap;
sensitivity to model parameters ($\kappa$, $\sigma(t)$, correlation) and to the
calibration set; stability of Greeks; benchmark against an independent model.
\iqlookfor{a structured validation plan.}
\end{solution}

\begin{solution}{iq:rc:bermudans-and-callables:6}
With a standard error of about two basis points, a two-basis-point move is
noise, not a bug. Report the standard error; use more paths, antithetics or
control variates (the Europeans have closed forms), fix the seed and the path
set for Greeks (common random numbers), and compare with the tree.
\iqlookfor{standard errors and variance reduction.}
\end{solution}
